\chapter{Lecture 3: Regular Expressions, Permissions, Variables and Scripts}
\begin{center}
\textit{January 11, 2018}
\end{center}

\section{Regular Expressions (continued)}
Last time: \code{egrep pattern file(s)}. For example
\begin{lstlisting}[style=bash]
egrep "(cs|CS)246" index.html
\end{lstlisting}
This is \emph{not} the same as \code{"(c|C)(s|S)246"} (which also matches
\code{cS246} and \code{Cs246}); that one can also be written
\code{"[cC][sS]246"}.

\begin{center}
\begin{tabular}{@{}lp{0.68\textwidth}@{}}
\toprule
Pattern & Matches \\
\midrule
\code{[abc]} & one character from the set, i.e.\ \code{(a|b|c)} \\
\code{[\^{}abc]} & one character \emph{not} from the set \\
\code{[a-z]} & one character in the range \\
\code{?} & 0 or 1 occurrences of the preceding expression \\
\code{*} & 0 or more occurrences of the preceding expression \\
\code{+} & 1 or more occurrences of the preceding expression \\
\code{.} & any one character \\
\code{.*} & any number of any characters \\
\code{\textbackslash.} & a literal \code{.} (backslash escapes a special symbol) \\
\code{\^{}} & forces the match to start at the beginning of the line \\
\code{\$} & forces the match to end at the end of the line \\
\bottomrule
\end{tabular}
\end{center}

Examples:
\begin{itemize}
  \item \code{"cs ?246"} matches \code{cs246} and \code{cs 246}.
  \item \code{"(cs)?246"} matches \code{cs246} or just \code{246}.
  \item \code{"\^{}cs246"}: lines starting with \code{cs246};
        \code{"\^{}cs246\$"}: lines containing exactly \code{cs246}.
\end{itemize}

\begin{example}
Print all words in \code{/usr/share/dict/words} that start with an \code{e} and
have length 5.
\begin{lstlisting}[style=bash]
egrep "^e....$" /usr/share/dict/words
\end{lstlisting}
\end{example}

\begin{example}
Lines of even length:
\begin{lstlisting}[style=bash]
egrep "^(..)*$" filename
\end{lstlisting}
\end{example}

\begin{important}[Anchors]
When a pattern must \emph{count} characters (exact length, exactly one
occurrence, \dots), anchor it with \code{\^{}} and \code{\$}. Otherwise
\code{egrep} happily matches a substring of the line.
\end{important}

\begin{example}
List the file names in the current directory that contain exactly one \code{a}.
\begin{lstlisting}[style=bash]
ls -a | egrep "^[^a]*a[^a]*$"
\end{lstlisting}
\end{example}

\section{File Permissions}
\code{ls -l} gives a long listing:
\begin{lstlisting}[style=bash]
-rwxr-xr--  1 nanaeem  staff  2048 Jan 11 10:00 script.sh
\end{lstlisting}
The columns are: type and permissions, number of links, owner, group, size,
last-modified date, and file name. The first character is \code{-} for an
ordinary file and \code{d} for a directory. The next nine characters are three
groups of three:

\begin{center}
\begin{tabular}{@{}ccc@{}}
\toprule
\code{rwx} & \code{r-x} & \code{r--} \\
user (owner) & group & other \\
\midrule
the file's owner & users in the file's group & everyone else \\
\bottomrule
\end{tabular}
\end{center}

\begin{center}
\begin{tabular}{@{}clp{0.42\textwidth}@{}}
\toprule
Bit & On an ordinary file & On a directory \\
\midrule
\code{r} & read the file & list the directory's contents (e.g.\ globbing) \\
\code{w} & modify the file & add/remove files in the directory \\
\code{x} & run the file as a program & enter the directory (\code{cd}) \\
\bottomrule
\end{tabular}
\end{center}

Only the owner can change a file's permissions, using \code{chmod mode file(s)}.
A mode consists of
\begin{enumerate}
  \item an ownership class: \code{u}ser, \code{g}roup, \code{o}ther, \code{a}ll;
  \item an operator: \code{+} (add), \code{-} (remove), \code{=} (set exactly);
  \item permissions: \code{r}, \code{w}, \code{x}.
\end{enumerate}
\begin{lstlisting}[style=bash]
chmod a+r final.pdf     # give everyone read permission
chmod g-rw final.pdf    # remove read and write from the group
chmod o=rx final.pdf    # set others to exactly read/execute
chmod u+x myscript      # make a script executable by me
\end{lstlisting}

\section{Shell Variables}
\begin{lstlisting}[style=bash]
x=1          # no spaces around = !  variables store strings
echo $x
echo ${x}    # recommended
\end{lstlisting}
When you \emph{write} a variable, no \code{\$}; when you \emph{read} it, use
\code{\$}. It is good practice to use curly braces around variable names:
\code{\$\{x\}abc} is unambiguous, \code{\$xabc} refers to a variable named
\code{xabc}.

\begin{itemize}
  \item Double quotes allow variable expansion; single quotes suppress it:
        \code{echo '\$\{x\}'} prints \code{\$\{x\}}.
  \item Variables last for the current session only. To set something
        permanently, put it in \code{\textasciitilde/.bash\_profile} (or
        \code{.bashrc}).
\end{itemize}

\subsection*{PATH}
\code{echo \$\{PATH\}} prints a list of directories separated by colons. When you
type a command, the shell searches these directories (in order) for a program
with that name.
\begin{lstlisting}[style=bash]
PATH=mypath:${PATH}    # search mypath first
PATH=${PATH}:mypath    # search mypath last
\end{lstlisting}

\section{Shell Scripts}
\begin{definition}
A \textbf{shell script} is a text file containing a sequence of shell commands
that can be executed as a program.
\end{definition}
\begin{itemize}
  \item The first line should be \code{\#!/bin/bash}, the \emph{shebang} (``hash
        bang'') line. It tells the OS to run the script with bash.
  \item Give the script executable permission: \code{chmod u+x script}.
  \item Run it with \code{./script} (the current directory is usually not in
        \code{PATH}).
\end{itemize}

\begin{lstlisting}[style=bash]
#!/bin/bash
date
whoami
pwd
\end{lstlisting}

\subsection{Arguments to a Script}
\begin{lstlisting}[style=bash]
./script arg1 arg2 ...
#  $0     $1   $2
\end{lstlisting}
\code{\$0} is the script's name, \code{\$1}, \code{\$2}, \dots{} are its arguments,
and \code{\$\#} is the number of arguments.

\begin{example}
Is a given word a valid word in the dictionary?
\begin{lstlisting}[style=bash]
#!/bin/bash
egrep "^$1$" /usr/share/dict/words
\end{lstlisting}
\end{example}

\subsection{Status Codes}
Every Linux command sets a \textbf{status code}, stored in \code{\$?}:
\code{0} means success, non-zero means failure. For \code{egrep}, success means
``found a match''.

\begin{aside}
To throw away output we don't care about, send it to \code{/dev/null}, a special
file that discards everything written to it.
\end{aside}

\subsection{Evaluating Conditions}
Conditions are evaluated with the \code{test} program, which can also be written
as \code{[ ... ]}. The spaces inside the brackets are mandatory, because
\code{[} is literally the name of a program and the rest are its arguments.

\begin{example}
A good password is a word that is \emph{not} in the dictionary. Is a given word a
good password?
\begin{lstlisting}[style=bash]
#!/bin/bash
# The answer is the script's own status code.
egrep "^$1$" /usr/share/dict/words > /dev/null
[ $? -ne 0 ]
\end{lstlisting}
The status of a script is the status of its last command, so the script
succeeds exactly when \code{egrep} failed to find the word.
\end{example}

\code{\#} starts a comment. \code{[ -e basic ]} checks whether a file named
\code{basic} exists.

\section*{Summary}
\begin{itemize}
  \item Regex: \code{[...]}, \code{[\^{}...]}, \code{?}, \code{*}, \code{+},
        \code{.}, \code{|}, grouping, and the anchors \code{\^{}} and \code{\$}.
  \item \code{ls -l} shows \code{rwx} permissions for user/group/other;
        \code{chmod} changes them.
  \item Variables: \code{x=1}, read with \code{\$\{x\}}; \code{PATH} is where
        commands are found.
  \item Scripts start with \code{\#!/bin/bash}, receive \code{\$1, \$2, ...}, and
        communicate success through status codes; \code{[ ]} tests conditions.
\end{itemize}
